% gerado por scripts/build_paper_assets.py -- nao editar a mao
\begin{table}[htbp]
\centering
\small
\caption{Individual models and model combinations on the same evaluation set, averaged
over the eight cities.}
\label{tab:ensembles}
\begin{tabular}{lrrrc}
\toprule
Approach & sMAPE (\%) & MAE & RMSE & Cities ranked first \\
\midrule
TimesFM & 74.9 & 3,082 & 10,862 & 5/8 \\
Median ensemble$\dagger$ & 79.0 & 3,392 & 10,908 & 1/8 \\
CatBoost & 83.4 & 3,691 & 11,165 & 1/8 \\
Inverse-error weighted ensemble$\dagger$ & 83.8 & 3,531 & 10,809 & 0/8 \\
XGBoost & 84.5 & 4,230 & 12,511 & 1/8 \\
Random Forest & 86.9 & 4,224 & 11,827 & 0/8 \\
Mean ensemble$\dagger$ & 87.6 & 3,787 & 10,979 & 0/8 \\
Stacked ensemble (Ridge)$\dagger$ & 94.8 & 3,526 & 10,532 & 0/8 \\
Ridge & 100.1 & 3,696 & 10,856 & 0/8 \\
SARIMAX & 108.6 & 4,070 & 11,531 & 0/8 \\
LightGBM & 116.2 & 5,435 & 11,874 & 0/8 \\
\bottomrule
\end{tabular}

\vspace{0.5em}
\begin{minipage}{\textwidth}
\footnotesize
$\dagger$ marks a combination of the seven individual models. Weights and the stacking
meta-model are fitted only on folds earlier than the one being predicted, so the
combination never sees the fold it is evaluated on. Individual models and combinations are
scored on exactly the same rows: the first 20 folds are withheld as a training window for
the combiners, leaving 1{,}188 predictions per approach per city. Comparing a combination
evaluated on 99 folds against a single model evaluated on 119 would favour the combination,
because the withheld folds are those with the shortest training window.
\end{minipage}
\end{table}
